% Paper 4, §12 — Cross-method convergence and divergence
% Thesis: across §3-§11 we have five methods (lens, probes, SAEs,
% activation patching, circuit tracing). This is the paper's central
% synthesis. When multiple methods converge on the same picture (IOI
% on GPT-2 small; factual recall on GPT-2 XL), the joint claim is
% load-bearing in a way no single method's claim is. When methods
% diverge — truth direction across datasets, SAE feature inventories
% across widths, lens cross-model unreliability — the divergence is
% itself diagnostic: either at least one method's assumptions have
% been falsified, or the methods are answering different questions
% and the apparent conflict is principled. Forward-link to §13 (the
% scaling/coverage map) and §14 (consolidated contradictions).

\section{Cross-method convergence and divergence}
\label{sec:cross-method}

The five methods of \S\S\ref{sec:logit-lens-family}--\ref{sec:circuits}
make different commitments. Lenses are correlational and project the
residual stream through the model's own unembedding. Probes are
correlational and project it through a trained classifier. Sparse
autoencoders are unsupervised and project it into a learned
$M\!\gg\!D$-dimensional sparse basis. Activation patching is causal at
the component level. Circuit tracing is causal at the feature level.
The methodological question this section asks is: when methods
\emph{should} agree --- because the underlying mechanism is the same
and they answer the same question about it --- what do their points
of convergence and divergence tell us about the model and about the
methods?

The answer is two-sided. \S\S\ref{sec:cross-ioi}--\ref{sec:cross-rome}
take up the convergence cases (IOI on GPT-2 small; factual recall
on GPT-2 XL), where three or four methods triangulate on the same
picture and the joint claim is more robust than any single-method
claim. \S\S\ref{sec:cross-truth}--\ref{sec:cross-lens} take up the
divergence cases (truth direction across datasets; SAE feature
canonicality; lens cross-model unreliability), where the methods
disagree and the disagreement is itself diagnostic.
\S\ref{sec:cross-principle} states the closing argument and
forward-links to \S\ref{sec:scaling-snapshot} and
\S\ref{sec:contradictions}.

\subsection{Convergence: the IOI circuit on GPT-2 small}
\label{sec:cross-ioi}

The most-replicated cross-method picture in the field is
indirect-object identification on GPT-2 small. The substrate is
$\sim 124$M parameters; the task is the canonical Wang--prompt
template ``\textsc{When Mary and John went to the store, John gave a
drink to \_\_}''; the target is the indirect-object name
\citep{wang2022-ioi}\footnote{KB excerpt:
\texttt{kb/excerpts/wang2022-ioi.md\#sec-2-task-and-circuit}.}. Three
methodologically independent readouts converge on the same circuit
picture.

\citet{wang2022-ioi} identify a circuit of $26$ attention heads in
$7$ functional classes (Duplicate Token, Previous Token,
S-Inhibition, Name Mover, Negative Name Mover, Backup Name Mover,
Induction Heads), $1.1\%$ of all (head, token-position) pairs, by
manual path-patching. \citet{conmy2023-acdc} recover a substantially
overlapping subgraph by automating the discovery: the ACDC
algorithm of \S\ref{sec:circuit-acdc} iterates over
computational-graph edges, KL-thresholds the indirect effect of
each, and prunes below threshold; the surviving subgraph on GPT-2
small recovers the Wang circuit's load-bearing edges without
human-supplied prior structure. A third method enters from the
lens side: the tuned-lens prediction trajectory of
\S\ref{sec:lens-trajectory} on IOI prompts shows the
indirect-object name's probability rising in the mid-late layer
band of GPT-2 small, consistent with the Name-Mover-and-S-Inhibition
commit identified by Wang et al.\ at the head level. SAE-feature
analysis on the same prompts identifies features matching the QK
and OV roles of the discovered heads
\deepencite{IOI SAE feature analysis: Bricken-style features matching
QK/OV roles, citation pending — Neuronpedia replication, not in
papers.json}.

\begin{intuition}
Three methods, three different commitments (causal-component,
causal-edge, correlational-readout), one converging picture: IOI on
GPT-2 small uses a sparse subset of attention heads in 7 functional
classes, distributed across mid-late layers, and the predicted
indirect-object probability emerges in exactly the layer band the
heads operate in. Returning to canonical form: the intersection of
patching's component-level localization, ACDC's edge-level discovery,
and the lens's prediction trajectory is the circuit
$C$ of \S\ref{sec:substrate-circuit}, satisfying faithfulness,
completeness, and minimality at the GPT-2-small scale.
\end{intuition}

\subsection{Convergence: factual recall on GPT-2 XL}
\label{sec:cross-rome}

The second positive cross-method case is factual recall on GPT-2 XL,
where four methods converge on the same localization. The substrate
is $1.5$B parameters; the task is the
$(s, r, o)$-triple cloze of \S\ref{sec:rome-tracing} (e.g., ``The
Eiffel Tower is in \_\_''); the target is the object completion.

\citet{meng2022-rome} run causal tracing across (component, token,
layer) cells and report two effects: the average indirect effect
peaks at MLP layers around layer 15--20 at the last subject token,
and the AIE for MLP perturbations ($6.6\%$) substantially exceeds
the AIE for attention perturbations ($1.6\%$) at that locus
\citep[\S 2.3]{meng2022-rome}\footnote{KB excerpt:
\texttt{kb/excerpts/meng2022-rome.md\#sec-2-effects}.}. The first
method --- activation patching with Gaussian noise as the corruption
--- localizes the mediating component class to mid-layer MLPs at the
last subject token.

The second method enters as the weight-level intervention. Following
the localized factual association hypothesis,
\citeauthor{meng2022-rome} construct a rank-one MLP-weight edit at
the localized layer that inserts a new $(s, r, o)$ triple under a
key--value constraint
$(W + \Delta W)\,\mathbf{k}_* = \mathbf{v}_*$ minimizing
$\|\Delta W \mathbf{k}_*\|^2$ \citep[\S 3.2]{meng2022-rome}. The
edit inserts new facts with $99.8\%$ efficacy and $88\%$ paraphrase
generalization on zsRE. The patching localization survives the
strongest possible test: an edit at the localized site succeeds at
inserting facts; an edit elsewhere does not. (See
\S\ref{sec:rome-hase-critique} for the Hase et al.\ 2023 ``does
localization inform editing'' caveat
\deepencite{hase2023-localization-not-editing §1, citation-pending}:
the \emph{precise} layer-of-edit is not always the AIE peak, but
the locus is robust.)

The third method enters from the lens side. The tuned-lens
prediction trajectory of \S\ref{sec:lens-trajectory} on factual
recall prompts shows the predicted-object token's logit rising
across exactly the localized layer band: the object is not yet
predicted at early layers, becomes predictable across the mid-layer
MLPs, and is committed by the late layers. The lens trajectory and
the patching heatmap are reading the same phenomenon through two
different windows.

The fourth method enters from the SAE side. SAE features trained on
the residual stream and MLP outputs of GPT-2 XL recover
subject-entity features at exactly the layer band patching
localizes \deepencite{Templeton-style SAE feature analysis on GPT-2 XL
factual-recall prompts; partial replication of
\citet{templeton2024-scaling-monosemanticity} methodology on a
smaller open model, citation pending}. The features are not fully
catalogued at the GPT-2 XL scale, but the structural pattern ---
subject-entity features are sparse in the localized band and not
elsewhere --- holds.

Four methods, one converging picture: factual recall in GPT-2 XL is
mediated by mid-layer MLPs at the last subject token; the effect
can be localized causally (patching), validated intervention-ally
(rank-one edit), seen correlationally (lens trajectory), and
decomposed at the feature level (SAE features). The joint claim is
more robust than any single-method claim because each method's
confound is different and the four do not align: a lens-only
artifact would not produce a successful edit; a patching-only
artifact (e.g., Gaussian-noise off-distribution effects
\citep{zhang2023-apatching}) would not align with SAE-discovered
features that are training-distribution by construction.

\subsection{Divergence: truth direction across datasets}
\label{sec:cross-truth}

The negative cases begin with the most-studied probing target. The
\citet{marks-tegmark-2023-truth} mass-mean direction on
LLaMA-2-13B, $\mathbf{d}_{\mathrm{MM}} = \boldsymbol{\mu}_+ -
\boldsymbol{\mu}_-$, is the canonical truth probe of
\S\ref{sec:mass-mean}; the same \S 4 of the paper documents
``stark misalignment'' across datasets \citep[\S 4]{marks-tegmark-2023-truth}\footnote{KB
excerpt: \texttt{kb/excerpts/marks-tegmark-2023-truth.md\#sec-4-misalignment}.}.
The truth direction recovered from the \texttt{cities} dataset
(``The capital of France is Paris.'') and the truth direction
recovered from \texttt{neg\_cities} (``The capital of France is
not Berlin.'') are antipodal in early layers, orthogonal in
mid-layers, and aligned only in late layers (around layers
12--15). Whether the recovered direction is a property of the
\emph{model}'s truth representation or of the \emph{dataset}'s
truth-pattern is layer-dependent.

A second method on the same target gives a different answer.
Contrast-Consistent Search (CCS)
\deepencite{burns2022-ccs §3, citation-pending} finds an
unsupervised direction satisfying a probabilistic-consistency loss
on contrast pairs; \citet[\S 1 related work]{marks-tegmark-2023-truth}\footnote{KB
excerpt: \texttt{kb/excerpts/marks-tegmark-2023-truth.md\#sec-1-related}.}
documents that the recovered direction is not the mass-mean
direction, and CCS suffers documented failures under negation that
mass-mean does not share. A third method on the same target enters from the SAE side.
SAE-discovered ``truthfulness''-related features
\citep[\S 4]{templeton2024-scaling-monosemanticity}\footnote{KB
excerpt: \texttt{kb/excerpts/templeton2024-scaling-monosemanticity.md\#sec-4}
(HTML, deepen-cite pending).} \deepencite{templeton2024-scaling-monosemanticity §4
truthfulness features, HTML-only quotes pending Phase 2 transcription}
overlap partially but not fully with mass-mean directions: the
cosine between the probe direction and the SAE feature decoder
column on matched concepts is, to the best of our knowledge, not
yet systematically reported, and the few hand-spot-check writeups
report partial alignment rather than identity.

\begin{contradiction}
\textbf{There may not be a single truth direction.} Mass-mean
probing on LLaMA-2-13B finds dataset-dependent directions that align
only in late layers \citep[\S 4]{marks-tegmark-2023-truth}; CCS
finds a different direction that fails under negation; SAE-discovered
truthfulness-related features overlap partially but not fully with
either. Three methods, no single converging direction. The
\citet{vennemeyer2025-sycophancy-not-one-thing} multi-component
result on sycophancy generalizes the pattern: behaviorally-singular
properties may be representationally plural, and the field's
single-axis truth-probe assumption may be the wrong shape of model
for the underlying representation. The strongest form of the open
question --- ``does the model represent truth along a single linear
axis, or along a multi-axis subspace whose components depend on
prompt format?'' --- is unresolved as of 2026, and the experiment
that would close it is a cross-method, cross-prompt-format,
cross-model-family study reporting probe direction, SAE feature
direction, and patching-validated causal direction in matched
coordinates.
\end{contradiction}

This is principled divergence in part. A probe asks ``what
direction linearly decodes property $y$ from this dataset's
activations?''; an SAE asks ``what sparse basis decomposes
activations under reconstruction-and-sparsity loss?''. Different
answers need not falsify each other --- but the divergence becomes
load-bearing once the field makes single-axis truth-probe claims
(``ablating $\mathbf{d}_y$ flips the model's TRUE/FALSE behavior''):
if there is no single $\mathbf{d}_y$, the claim is
dataset-conditional. The divergence is the diagnostic.

\subsection{Divergence: SAE feature canonicality}
\label{sec:cross-sae}

A second negative case is the SAE feature inventory itself.
\citet{leask2025-sae-not-canonical} show that two independent SAE
training runs at different widths $M$ may produce different
feature inventories on the same activations
\citep[\S 1]{leask2025-sae-not-canonical}\footnote{KB excerpt:
\texttt{kb/excerpts/leask2025-sae-not-canonical.md\#sec-1-canonical}.}:
meta-SAEs decompose larger-SAE latents into combinations of
smaller-SAE latents (non-atomicity), and SAE stitching reveals
novel latents present in larger SAEs but absent from smaller ones
(incompleteness). The
\citet{bricken2023-monosemanticity}/\citet{templeton2024-scaling-monosemanticity}
``true features'' framing assumes a sufficiently large SAE
recovers the model's underlying units of computation modulo
training noise; the \citeauthor{leask2025-sae-not-canonical}
critique argues that ``sufficiently large'' is not well-defined.
\S\ref{sec:sae-canonicality} treats the contradiction at depth;
the cross-method consequence is that any claim grounding in SAE
features (``the SAE feature for truth has cosine $\rho$ with the
probe direction'') is a claim about \emph{this} SAE's inventory,
not about a model-canonical basis. The pragmatic 2026 position ---
training SAEs at multiple widths and reporting which claims survive
across widths --- reduces the single-run-artifact risk; the open
theoretical question (a basis-finding criterion under which
independent runs converge modulo a known equivalence) is
documented in \S\ref{sec:canonicality-frontier}.

\subsection{Divergence: lens cross-model unreliability}
\label{sec:cross-lens}

The third negative case is internal to the lens family. The logit
lens of \citet{nostalgebraist2020-logit-lens} works on GPT-2 ---
the trajectory $\{\mathrm{LogitLens}(\mathbf{h}_\ell)\}_{\ell=0}^L$
converges smoothly toward the final-layer prediction --- but on
BLOOM and GPT-Neo of equivalent scale fails: intermediate
distributions are systematically biased and the no-training
baseline is no longer a useful diagnostic
\citep[\S 2.2]{belrose2023-tuned-lens}\footnote{KB excerpt:
\texttt{kb/excerpts/belrose2023-tuned-lens.md\#sec-2-failures}.}.
\citet{belrose2023-tuned-lens} fix the bias with the per-layer
affine translator
$\mathrm{TunedLens}_\ell(\mathbf{h}_\ell) = \mathrm{LogitLens}(A_\ell\,
\mathbf{h}_\ell + \mathbf{b}_\ell)$, at the cost of
$L\!\cdot\!(D^2 + D)$ trained parameters.
\citet{chuang2023-dola}'s DoLa contrastive decoding works on the
LLaMA-7B/13B/33B family; whether it generalizes to Llama 3 / 4,
Qwen 2.5 / 3, DeepSeek-V3, and Gemma 3 has not been systematically
reproduced, and the \texttt{logitlens4llms-2025} extension is the
partial cover for the 2025--2026 frontier. The cross-method
consequence: any synthesis using a lens trajectory (the IOI and
factual-recall convergence cases above) is implicitly
model-conditional, and the four-method factual-recall convergence
of \S\ref{sec:cross-rome} is therefore a GPT-2-XL-specific result
until a tuned lens for the frontier-2026 architecture exists.

\subsection{When divergence is principled, and when it is falsification}
\label{sec:cross-principle}

The closing argument is structural. Two methods can disagree for
two reasons: they answer the same question with different
commitments and one is wrong (\emph{falsification divergence}, and
the diagnostic action is to identify which method's assumption
breaks); or they answer different questions whose answers happen
to disagree on a target (\emph{principled divergence}, and the
diagnostic action is to be careful about which method's answer one
cites for which claim). The convergence cases above are the regime
where methods are answering compatible questions and the answers
cross-validate. The truth-direction divergence is on the boundary:
probes and SAE features ask different questions, but the
single-axis truth claim assumes the questions have a common answer,
and the divergence falsifies that assumption (or exposes that the
representation is multi-axis). The SAE-canonicality and lens
cross-model divergences are internal to one method family: two
runs of the same procedure return different answers, and the
divergence falsifies a stability assumption in each case.

\begin{intuition}
Method-pluralism is principled because the methods commit
differently and answer different questions. But the convergence
cases above show that, for sufficiently constrained behaviors on
sufficiently small models, the methods \emph{can} be made to
answer the same question, and when they do, agreement is the
load-bearing form of evidence. Returning to canonical form: the
joint posterior over circuits, given evidence from
$\{\text{patching}, \text{lens}, \text{SAEs}, \text{ROME-edit}\}$,
is sharply peaked when all four methods agree (the IOI and
factual-recall cases) and diffusely supported when they do not
(the truth-direction case). The next paper-level methodological
move is not better single methods --- it is principled
cross-method protocols that fix what each method is asked to
answer, so that disagreement is interpretable as falsification
rather than as a question mismatch.
\end{intuition}

This synthesis feeds two downstream sections.
\S\ref{sec:scaling-snapshot} consumes the convergence cases as the
regime where every method has scaled enough to participate, and
the divergence cases as the regime where single-method claims are
still the only evidence on production-scale models.
\S\ref{sec:contradictions} consolidates the divergence cases of
\S\S\ref{sec:cross-truth}--\ref{sec:cross-lens} alongside the
\S\S\ref{sec:circuit-disagreement}, \ref{sec:canonicality-frontier},
and \ref{sec:rome-hase-critique} contradictions surfaced earlier
in the paper, and pairs each with the experiment that would close
the gap. The thesis holds: ``\emph{partially read}'' is the right
phrase because on the well-constrained behaviors where the methods'
angles coincide (IOI, factual recall) the model is read tightly,
while on the open behaviors where they do not (truth, sycophancy,
the SAE feature inventory) the field has identified the
disagreement and is working on the resolution.
